% ===================================================================
%  ТАБЛИЦЯ № 6 — Дім усередині. — Меблі. — Їжа. — Освітлення.
%  Traduction 2026 d'apres texto/io, controlee sur texto/fr, texto/en
%  et texto/es. Consigne : voir texto/uk/10-tablycia-01.tex.
%
%  Deux notes, toutes deux marquees « (*) », et deux appels : celui de
%  « babo » au bloc c1-12-1, celui de « lambrequino » au bloc c2-01-1.
%  L'ordre les departage, comme dans les autres colonnes.
%
%  LA FEMME DE CHAMBRE. L'ido du bloc c1-06-1 porte « (6) » la ou le
%  francais porte « (16) », et c'est le francais qui a raison.
%  L'ukrainien suit le francais et ecrit (16), comme toutes les autres
%  colonnes ; gravuri/korekti.json redresse le renvoi ido dans la page
%  de lecture, et renvoji.py declare l'ecart.
%
%  LE PLUS GROS LEXIQUE DE LA COLONNE — cinq pieces, deux cent
%  cinquante objets — et le plus riche en pieges. La cuisine surtout :
%  la passoire est un « друшляк », la rape une « тертка », la louche
%  un « ополоник », l'evier une « зливальниця », le hachoir un
%  « сікач », la lessiveuse une « балія ». Six ustensiles, six mots
%  que le russe ne dit pas ainsi.
% ===================================================================

%%K t06-apar-1 apar
\VUsaut{6.0mm}
\VUtitre{15.0pt}{60}{ТАБЛИЦЯ № 6}
\VUsaut{1.4mm}
\VUtitre{11.4pt}{0}{Дім усередині, Меблі, Їжа,}
\VUsaut{0.4mm}
\VUtitre{11.4pt}{0}{Освітлення.}
\VUsaut{2.2mm}

%%K t06-apar-2 apar
Дім готовий. Іванів дядько перебрався до свого нового житла, розташування
якого ми вже знаємо. Подивімося тепер, як він його обставив. Спершу ми
відвідаємо:

%%K t06-tit-1 sub
\VUpk{10.2pt}{40}{Спальню.}
\VUsaut{3.0mm}

%%K t06-c1-01-1 p
1. --- Ми прийшли невдалої хвилини, бо служниця прибирає \VUgras{кімнату}.

%%K t06-c1-02-1 p
2. --- На \VUgras{умивальнику} \textsuperscript{(1)} лежать там і сям різні речі,
якими миються, чешуться і взагалі опоряджаються. Це \VUgras{мідниця}
\textsuperscript{(2)} і \VUgras{глек} \textsuperscript{(3)}, \VUgras{щітка до
волосся} \textsuperscript{(4)}, \VUgras{гребінець} \textsuperscript{(5)},
\VUgras{мило} \textsuperscript{(6)} і \VUgras{губка} \textsuperscript{(7)},
нарешті \VUgras{пляшечка} \textsuperscript{(8)} з полосканням до рота.

%%K t06-c1-03-1 p
3. --- На підлозі, перед умивальником, стоїть \VUgras{відро}
\textsuperscript{(9)} на брудну воду.

%%K t06-c1-04-1 p
4. --- У \VUgras{передпокої} \textsuperscript{(10)} видно кілька \VUgras{вішалок}
\textsuperscript{(13)} і дві \VUgras{парасольки} \textsuperscript{(11)} у
\VUgras{підставці} \textsuperscript{(12)}.

%%K t06-c1-05-1 p
5. --- По другий бік дверей стоїть \VUgras{дзеркальна шафа}
\textsuperscript{(14)}, у якій зберігається \VUgras{білизна}
\textsuperscript{(15)}.

%%K t06-c1-06-1 p
6. --- Скористаймося з того, що \VUgras{покоївка} \textsuperscript{(16)} стеле
\VUgras{постіль} \textsuperscript{(17)}, і розглянемо її частини.
\VUgras{Ліжко} \textsuperscript{(20)} несе добротну \VUgras{пружинну сітку}
\textsuperscript{(21)}, на яку Юстина зараз покладе вовняний \VUgras{матрац}
\textsuperscript{(22)}. Потім вона візьме зі стільців два \VUgras{простирадла}
\textsuperscript{(23)} і \VUgras{ковдру} \textsuperscript{(24)}. Тоді покладе в
головах \VUgras{валик} \textsuperscript{(25)} і \VUgras{подушку}
\textsuperscript{(26)}, які лежать разом із \VUgras{периною}
\textsuperscript{(27)} на \VUgras{приліжковому килимку} \textsuperscript{(32)}.
Мітелкою вона обмахне \VUgras{завіси} \textsuperscript{(19)} і
\VUgras{балдахін} \textsuperscript{(18)}, і частину роботи зроблено.

%%K t06-c1-07-1 p
7. --- Потім їй доведеться трохи прибрати \VUgras{нічний столик}
\textsuperscript{(28)}, накрутити \VUgras{будильник} \textsuperscript{(29)},
приготувати \VUgras{нічник} \textsuperscript{(30)} і винести \VUgras{свічку}
\textsuperscript{(31)}, щоб вичистити свічник.

%%K t06-tit-2 sub
\VUsaut{5.0mm}
\VUpk{10.2pt}{40}{Кошик з рукоділлям і камінне начиння.}
\VUsaut{3.0mm}

%%K t06-c1-08-1 p
8. --- \VUgras{Кошик з рукоділлям} \textsuperscript{(34)} біля \VUgras{табурета}
\textsuperscript{(33)} теж дуже потребує прибирання. Доведеться скласти
\VUgras{вишивку} \textsuperscript{(35)}, поховати в \VUgras{робітний столик}
\textsuperscript{(38)} \VUgras{наперсток}, \VUgras{нитки} \textsuperscript{(36)},
\VUgras{голки} \textsuperscript{(37)} і \VUgras{ножиці} \textsuperscript{(39)} і
замкнути накривку столика \VUgras{ключем} \textsuperscript{(40)}.

%%K t06-c1-09-1 p
9. --- На \VUgras{круглому столику} \textsuperscript{(41)} посеред кімнати Юстина
змінить воду в \VUgras{карафці} \textsuperscript{(42)}. Вона покладе
\VUgras{цукор} у \VUgras{цукерницю} \textsuperscript{(43)}, вичистить
\VUgras{щипчики до цукру} \textsuperscript{(44)} і винесе \VUgras{щітку до одягу}
\textsuperscript{(46)}.

%%K t06-c1-10-1 p
10. --- Правий бік кімнати дасть їй менше роботи, бо \VUgras{канапка}
\textsuperscript{(47)} і її \VUgras{подушка} \textsuperscript{(48)} на своїх
місцях, а що не холодно, то нічого не чіпали і в начинні \VUgras{каміна}
\textsuperscript{(49)}. \VUgras{Совок} \textsuperscript{(50)}, \VUgras{щипці}
\textsuperscript{(51)}, \VUgras{тагани} \textsuperscript{(55)} і \VUgras{решітка
на попіл} \textsuperscript{(56)} чисті; \VUgras{камінний екран}
\textsuperscript{(54)} не зсунуто, а \VUgras{скриня на дрова}
\textsuperscript{(52)} повна \VUgras{дров} \textsuperscript{(53)}.

%%K t06-noto-1 noto
\VUnotes{143.2mm}{%
(*) Babo = маленька дитина; англ. \textit{baby}, фр. \textit{bébé}, іт.
\textit{bambino}.%
}

%%K t06-noto-2 noto
\VUnotes{143.2mm}{%
(*) Lambrequino = фр. \textit{lambrequin}, ісп. \textit{lambrequines}, порт.
\textit{lambrequins}, і навіть нім. і англ. \textit{lambrequins} --- тобто і
по-німецькому, і по-англійському, і по-французькому, і по-португальському, і
по-іспанському.%
}

%%K t06-c1-11-1 p
11. --- Їй ще треба буде обмахнути \VUgras{камінний годинник}
\textsuperscript{(57)}, \VUgras{канделябри} \textsuperscript{(58)} і
\VUgras{тацівки на дрібнички} \textsuperscript{(59)}, а потім протерти
\VUgras{дзеркало} \textsuperscript{(60)}.

%%K t06-c1-12-1 p
12. --- Що ж до \VUgras{колиски} \textsuperscript{(61)} малюка (бабо (*)), вона
вже готова.

%%K t06-tit-3 sub
\VUsaut{5.0mm}
\VUpk{10.2pt}{40}{Вітальню.}
\VUsaut{3.0mm}

%%K t06-c2-01-1 p
1. --- А ось і вітальня. Це дуже гарна \VUgras{кімната}. Камін у правому кутку
прикрашений зграбною \VUgras{статуеткою} \textsuperscript{(1)} і двома гарними
\VUgras{вазами} \textsuperscript{(3)}, а задрапований оксамитовим
\VUgras{ламбрекеном} \textsuperscript{(2)} (*).

%%K t06-c2-02-1 p
2. --- Щоб іскри з вогнища не підпалили килима, перед каміном поставлено мідний
\VUgras{екран} \textsuperscript{(4)}.

%%K t06-c2-03-1 p
3. --- Поруч стоїть розкрите зграбне дамське \VUgras{бюрко}
\textsuperscript{(5)}. Тут \VUgras{господиня дому} \textsuperscript{(23)} пише
свої листи.

%%K t06-c2-04-1 p
4. --- Вікно прибране \VUgras{тюлевими фіранками} \textsuperscript{(6)},
підхопленими \VUgras{підхватами} \textsuperscript{(8)} на \VUgras{гачках}
\textsuperscript{(7)}, і важкими матерчатими шторами з \VUgras{підзором}
\textsuperscript{(9)} згори.

%%K t06-c2-05-1 p
5. --- У ніші вікна \VUgras{жардиньєрка} \textsuperscript{(11)} тримає зелену
\VUgras{рослину} \textsuperscript{(10)}, чудову пальму, листя якої сягає майже до
\VUgras{гасової лампи} \textsuperscript{(12)}.

%%K t06-c2-06-1 p
6. --- \VUgras{Абажур} \textsuperscript{(13)} зробила Марія, чарівна
\VUgras{панночка} \textsuperscript{(14)} за \VUgras{роялем}
\textsuperscript{(15)}. Сидячи дуже прямо на \VUgras{табуреті}
\textsuperscript{(16)}, вона розучує п'єсу. Вона дуже вправна і дуже охайна, про
що каже її \VUgras{нотна шафка} \textsuperscript{(17)}: там повний лад.

%%K t06-tit-4 sub
\VUsaut{5.0mm}
\VUpk{10.2pt}{40}{Пустуна.}
\VUsaut{3.0mm}

%%K t06-c2-07-1 p
7. --- Її брат Павло шукає \VUgras{книжку} \textsuperscript{(19)} у
\VUgras{книжковій шафі} \textsuperscript{(18)}. Це \VUgras{хлопець}
\textsuperscript{(20)} непосидючий і зовсім нестерпний. Повиснувши на шафі, щоб
дістати надто високо поставлену книжку, він важить звалити \VUgras{погруддя}
\textsuperscript{(21)}, що стоїть нагорі.

%%K t06-c2-08-1 p
8. --- Він користається для цієї дурниці з того, що мати захопилася розмовою з
приятелькою, \VUgras{панією} \textsuperscript{(22)}, яка приїхала погостювати в
неї кілька днів. Мати сидить на \VUgras{канапі} \textsuperscript{(24)} біля
\VUgras{крісла} \textsuperscript{(25)}, а її приятелька на \VUgras{стільці}
\textsuperscript{(27)}, від якого нам видно тільки \VUgras{спинку}
\textsuperscript{(26)}.

%%K t06-c2-09-1 p
9. --- У великій \VUgras{рамі} \textsuperscript{(30)} на стіні ---
\VUgras{портрет} \textsuperscript{(29)} родича. В \VUgras{альбомі}
\textsuperscript{(32)}, що лежить на столі, зібрано \VUgras{світлини}
\textsuperscript{(33)} решти рідні. Діти щойно його гортали, бо \VUgras{стільці}
\textsuperscript{(31)} ще зсунуті навколо столу.

%%K t06-c2-10-1 p
10. --- Порцелянову \VUgras{китайську вазу} \textsuperscript{(34)} біля
\VUgras{ножа до паперу} \textsuperscript{(35)} подарували Марії на іменини, а
\VUgras{ширму} \textsuperscript{(36)}, що займає кут кімнати, привіз із Китаю її
дядько.

%%K t06-c2-11-1 p
11. --- Оживлена гарними \VUgras{картинами} \textsuperscript{(37)} на
\VUgras{стіні} \textsuperscript{(42)}, освітлена стельовою \VUgras{люстрою}
\textsuperscript{(38)}, чиї \VUgras{електричні лампочки} \textsuperscript{(39)}
так весело світять вечорами, ця вітальня здається дуже приємною. М'який
\VUgras{килим} \textsuperscript{(40)}, розстелений по підлозі, і такий зручний
\VUgras{електричний дзвінок} \textsuperscript{(41)} роблять її цілком затишною.

%%K t06-tit-5 sub
\VUsaut{3.6mm}
\VUpk{10.2pt}{40}{Їдальню.}
\VUsaut{3.0mm}

%%K t06-c3-01-1 p
1. --- Подзвонили на обід. Уся родина, крім Марії, зібралася навколо столу.
Їдальня приємно натоплена добірною кахляною \VUgras{піччю} \textsuperscript{(1)}
з тонким \VUgras{димарем} \textsuperscript{(2)}.

%%K t06-c3-02-1 p
2. --- У цій кімнаті небагато меблів. Уздовж стіни, над \VUgras{табуретом}
\textsuperscript{(4)}, висить ошатний \VUgras{умивальний фонтанчик}
\textsuperscript{(3)}. Поруч стоїть \VUgras{буфет} \textsuperscript{(5)}, на який
ставлять холодні страви, десерти і різний посуд, поки не потрібний.

%%K t06-c3-03-1 p
3. --- Так, на верхній полиці ми бачимо \VUgras{печену курку}
\textsuperscript{(7)}, \VUgras{рибу} \textsuperscript{(8)} і \VUgras{вазу}
\textsuperscript{(10)} з \VUgras{садовиною} \textsuperscript{(9)}. Нижче стоять
\VUgras{сир} \textsuperscript{(11)}, \VUgras{хліб} \textsuperscript{(12)} і
\VUgras{судок} \textsuperscript{(13)}, у якому є все, що треба, щоб заправити
салат: олія, оцет, \VUgras{сіль} \textsuperscript{(14)} і \VUgras{перець}
\textsuperscript{(15)}. Нарешті, на нижній полиці стоїть \VUgras{пиріг}
\textsuperscript{(17)} поруч із \VUgras{гірчичницею} \textsuperscript{(16)}.

%%K t06-c3-04-1 p
4. --- Годинник у їдальні, який зветься \VUgras{стінним}
\textsuperscript{(20)}, дуже незвичайної форми. Його вставлено в дерев'яну раму,
у якій є також \VUgras{барометр} \textsuperscript{(19)} і \VUgras{термометр}
\textsuperscript{(18)}.

%%K t06-tit-6 sub
\VUsaut{4.6mm}
\VUpk{10.2pt}{40}{Як подають обід.}
\VUsaut{3.0mm}

%%K t06-c3-05-1 p
5. --- Навкруги всієї кімнати стіни обшито \VUgras{дерев'яними панелями}
\textsuperscript{(21)} з вощеного дуба. Матерчата \VUgras{портьєра}
\textsuperscript{(22)} доволі добре затуляє двері, що ведуть на веранду. Туди
родина піде після обіду пити каву. \VUgras{Чашки} \textsuperscript{(24)} і
\VUgras{блюдця} \textsuperscript{(25)} уже розставлено на \VUgras{гірці}
\textsuperscript{(23)}.

%%K t06-c3-06-1 p
6. --- Над столом до стелі прикріплено \VUgras{висячу лампу}
\textsuperscript{(26)}; її дуже яскраве світло заливає вечорами всю їдальню.

%%K t06-c3-07-1 p
7. --- Погляньмо тепер на сам \VUgras{обідній стіл} \textsuperscript{(27)}. Коли
родина ввійшла, стіл був накритий. На \VUgras{скатертині} \textsuperscript{(28)}
поставили на кожне місце \VUgras{тарілку} \textsuperscript{(33)},
\VUgras{серветку} \textsuperscript{(29)}, \VUgras{ложку} \textsuperscript{(34)},
\VUgras{виделку} \textsuperscript{(35)}, \VUgras{ніж} \textsuperscript{(36)} і
\VUgras{склянку} \textsuperscript{(32)}. Були там і \VUgras{пляшки}
\textsuperscript{(30)} на вино і \VUgras{карафка} \textsuperscript{(31)} на воду.

%%K t06-c3-08-1 p
8. --- \VUgras{Слуга} \textsuperscript{(39)} спершу вніс \VUgras{супову миску}
\textsuperscript{(37)}, у якій ще парує суп. Тепер він подає першу
\VUgras{страву} \textsuperscript{(38)}, м'ясну. Потім він обнесе всіх тими, що ми
вже назвали; і нарешті, після \VUgras{десерту}, скористається з того, що
господарі перейшли на веранду, щоб прибрати зі столу і знову дати лад їдальні.

%%K t06-c3-08-2 p
\textit{Примітка.} --- \textit{Ужиткові слова про їжу подано тут лише в
загальних рисах. Список буде доповнено у двох таблицях «Готель» (№ 13) і
«Ринок» (№ 15).}

%%K t06-tit-7 sub
\VUsaut{4.6mm}
\VUpk{10.2pt}{40}{Кухню.}
\VUsaut{3.0mm}

%%K t06-c4-01-1 p
1. --- На кухні, куди ми зазираємо крізь прочинені двері, нас зустрічає
мальовничий безлад. Перед \VUgras{вогнищем} \textsuperscript{(1)}, де тріщить
\VUgras{вогонь} \textsuperscript{(3)}, встановлено \VUgras{рожновий станок}
\textsuperscript{(5)}. Прекрасна \VUgras{печеня} \textsuperscript{(7)} надіта на
\VUgras{рожен} \textsuperscript{(6)} і доходить.

%%K t06-c4-02-1 p
2. --- \VUgras{Міх} \textsuperscript{(8)} і \VUgras{совок на вугілля}
\textsuperscript{(9)}, якими щойно користувалися, зосталися на підлозі разом із
\VUgras{полінами} \textsuperscript{(10)}.

%%K t06-c4-03-1 p
3. --- На камінній полиці лад: \VUgras{ліхтар} \textsuperscript{(11)},
\VUgras{сірничниця} \textsuperscript{(13)} із \VUgras{сірниками}
\textsuperscript{(12)}, \VUgras{свічник} \textsuperscript{(14)} і \VUgras{нічний
свічник} \textsuperscript{(15)} зі своєю \VUgras{свічкою} \textsuperscript{(16)}
--- усе на своїх місцях. Зате велике \VUgras{відро на вугілля}
\textsuperscript{(18)}, повне \VUgras{вугілля} \textsuperscript{(17)}, по яке
\VUgras{куховарка} \textsuperscript{(23)} щойно спускалася в підвал, лишено
посеред кухні.

%%K t06-c4-04-1 p
4. --- Річ у тім, що бідній жінці треба поспішати, щоб сніданок наспів вчасно.
Зараз вона коло \VUgras{плити} \textsuperscript{(19)}, помішує \VUgras{підливу}
\textsuperscript{(24)}, якій не можна дати пригоріти.

%%K t06-c4-05-1 p
5. --- У \VUgras{духовці} \textsuperscript{(20)} печеться пиріг, який вона
зробила дітям на підвечірок. Над плитою висять \VUgras{ополоник}
\textsuperscript{(21)} і \VUgras{шумівка} \textsuperscript{(22)}.

%%K t06-tit-8 sub
\VUsaut{4.0mm}
\VUpk{10.2pt}{40}{Кухонний посуд.}
\VUsaut{3.0mm}

%%K t06-c4-06-1 p
6. --- \VUgras{Набір каструль} \textsuperscript{(25)}, що висить у глибині
кімнати, дуже чистий. Гарні мідні \VUgras{каструлі} \textsuperscript{(26)},
\VUgras{молочник} \textsuperscript{(27)}, \VUgras{друшляк} \textsuperscript{(28)}
і \VUgras{тертка} \textsuperscript{(29)} старанно вичищені.

%%K t06-c4-07-1 p
7. --- \VUgras{Сільниця} \textsuperscript{(30)} стоїть на маленькому поставці
поруч із \VUgras{чайником} \textsuperscript{(31)} і \VUgras{сулією на олію}
\textsuperscript{(32)}. У \VUgras{шухляді} \textsuperscript{(33)}, треба гадати,
лежать прибори і кухонні ножі.

%%K t06-c4-08-1 p
8. --- \VUgras{Мітла} \textsuperscript{(34)} і \VUgras{мітелка з пір'я}
\textsuperscript{(35)} стоять у кутку. Куховарка щойно користувалася
\VUgras{колодою} \textsuperscript{(36)}; вона лишила її коло вікна.

%%K t06-c4-09-1 p
9. --- \VUgras{Рушник} \textsuperscript{(37)} висить на стіні поруч із
\VUgras{лійкою} \textsuperscript{(38)}. У цій частині кухні тримають
\VUgras{рашпер} \textsuperscript{(39)}, \VUgras{сікач} \textsuperscript{(40)} і
\VUgras{обробну дошку} \textsuperscript{(41)}. Тоді, над сікачем, на
\VUgras{полиці} \textsuperscript{(42)}, стоять \VUgras{горщики}
\textsuperscript{(43)}, \VUgras{чавунець} \textsuperscript{(44)} зі своєю
\VUgras{накривкою} \textsuperscript{(45)} і \VUgras{казан} \textsuperscript{(46)}.
\VUgras{Сковорода} \textsuperscript{(47)} висить поруч.

%%K t06-c4-10-1 p
10. --- Тільки мідний \VUgras{кран} \textsuperscript{(48)} кидає відблиск у цей
кут. \VUgras{Зливальниця} \textsuperscript{(49)}, на якій стоїть великий
\VUgras{глек} \textsuperscript{(50)}, мабуть, дуже зручна зі своєю шафкою, де
тримають \VUgras{барильце з оцтом} \textsuperscript{(51)} з його
\VUgras{краником} \textsuperscript{(52)}, \VUgras{скриньку на сміття}
\textsuperscript{(53)} і \VUgras{брудний посуд} \textsuperscript{(54)}.

%%K t06-tit-9 sub
\VUsaut{4.0mm}
\VUpk{10.2pt}{40}{Що ще тримають на кухні.}
\VUsaut{3.0mm}

%%K t06-c4-11-1 p
11. --- Велика \VUgras{балія} \textsuperscript{(55)}, у якій миють
\VUgras{городину} \textsuperscript{(58)}, і чавунний \VUgras{казан}
\textsuperscript{(56)}, що стоїть на \VUgras{кахляній підлозі}
\textsuperscript{(57)}, теж, певно, належать до цієї шафки.

%%K t06-c4-12-1 p
12. --- На плити куховарці, звісно, не бракує, бо ось тут, на розі столу, стоїть
\VUgras{газовий пальник} \textsuperscript{(59)}, а на ньому гріється вода в
каструльці. \VUgras{Пара} \textsuperscript{(60)}, що виходить із неї, каже нам,
що вода, мабуть, кипить. Ця вода, певно, на каву, бо ось на другому кінці столу
стоять \VUgras{кавник} \textsuperscript{(64)} і \VUgras{млинок до кави}
\textsuperscript{(63)}.

%%K t06-c4-13-1 p
13. --- Пальник з'єднано з \VUgras{газовим ріжком} \textsuperscript{(62)} довгою
\VUgras{гумовою руркою} \textsuperscript{(61)}.

%%K t06-c4-14-1 p
14. --- \VUgras{Поденниця} \textsuperscript{(66)}, яка приходить помагати
куховарці, готує городину на обід. Коли її буде обчищено і вимито, вона покладе
її в \VUgras{каструлю} \textsuperscript{(65)} до тієї пори, поки не прийде час її
варити.

%%K t06-tit-10 sub
\VUsaut{4.0mm}
{\centering\fontsize{10.2pt}{10.2pt}\selectfont\textls[40]{\scshape Ванну.} \textit{(Кімната для купання.)}\par}
\VUsaut{3.0mm}

%%K t06-c5-01-1 p
1. --- Лишається зазирнути ще в одне місце, щоб знати головні кімнати дому:
подивитися на влаштування ванної. Це недовго, бо вона невелика, і ми швидко
побачимо, що у \VUgras{ванни} \textsuperscript{(1)} є добра \VUgras{колонка}
\textsuperscript{(2)}, що \VUgras{мильниця} \textsuperscript{(3)} у того, хто
купається, під рукою, а \VUgras{вішалка на рушники} \textsuperscript{(5)} має
полегшувати сушіння \VUgras{рушників} \textsuperscript{(4)}.

%%K t06-c5-02-1 p
2. --- У цієї кімнати стіни вкриті \VUgras{полив'яною плиткою}
\textsuperscript{(6)}, що дозволило поставити \VUgras{душ} \textsuperscript{(7)},
не боячись, що вода, яка бризкає при користуванні, їх зіпсує. Маленька цинкова
\VUgras{мідниця} \textsuperscript{(8)} для ножних ванн довершує обстановку.
\VUsaut{6.0mm}
\VUfilet{22mm}
